\begin{lemma}
Let \(M\) be a finite family of medium base vertices, let
\(w(x)\in\mathbb N\) be the number of incidences charged to \(x\), and let
\(E(B)\in\mathbb N\) be the number of distinct present base edges charged by a
subfamily \(B\subseteq M\).  Suppose \(A\ge0\), \(L>0\),
\[
A<\sum_{x\in M} w(x),
\]
each \(x\in M\) satisfies \(w(x)\le t_2\) and
\(\ell\le w(x)\), and every subfamily satisfies the charge inequality
\[
\sum_{x\in B} w(x)\le L\,E(B).
\]
Then there is a nonempty subfamily \(B\subseteq M\) such that
\[
A<\sum_{x\in B} w(x)\le A+t_2,\qquad
|B|\,\ell\le A+t_2,\qquad
E(B)>A/L.
\]
\end{lemma}
\begin{proof}
Choose, among all subfamilies \(C\subseteq M\) with
\(A<\sum_{x\in C}w(x)\), one of minimum cardinality, and call it \(B\).
Such a \(B\) exists because the full family \(M\) has sum larger than \(A\).
The empty family has sum \(0\le A\), so \(B\) is nonempty.

Fix \(b\in B\).  By the cardinal-minimal choice of \(B\), deleting \(b\)
cannot leave a subfamily whose sum is still larger than \(A\).  Hence
\[
\sum_{x\in B\setminus\{b\}}w(x)\le A.
\]
Adding back the single term \(w(b)\), and using \(w(b)\le t_2\), gives
\[
\sum_{x\in B}w(x)\le A+t_2.
\]
The lower bound \(\ell\le w(x)\) for each \(x\in B\) gives
\[
|B|\ell\le \sum_{x\in B}w(x)\le A+t_2.
\]
Finally, the charge inequality for this same \(B\) gives
\[
A<\sum_{x\in B}w(x)\le L E(B).
\]
Since \(L>0\), division by \(L\) yields \(E(B)>A/L\).
\end{proof}
